% Source: Wald GR, physical PDF p. 106
%         (printed p. 93).
% Coverage: Figure 5.2.
% Figures/tables/footnotes: Figure 5.2; no tables or footnotes.
% Status: source-reviewed and compile-clean.
% Uncertainties: none.
\begingroup
\renewcommand{\thefigure}{5.2}
\begin{figure}[htbp]
  \centering
  \begin{tikzpicture}[>=Stealth, line cap=round, line join=round, scale=.90]
    \foreach \x/\bend in {-2.3/-16,-1.55/-10,-.78/-5,0/0,.78/5,1.55/10,2.3/16} {
      \draw[thick] (\x,-1.30) .. controls (\x+\bend/35,.15) and
        (\x-\bend/35,1.20) .. (\x,2.85);
    }
    \fill (0,.62) circle (1.5pt) node[below left=2pt] {\(p\)};
    \draw[->, thick] (0,.62) -- (.05,1.52) node[left] {\(u^a\)};
    \draw[->, thick] (0,.62) -- (.62,1.08) node[above right] {\(s_1^a\)};
    \draw[->, thick] (0,.62) -- (1.02,.62) node[below right] {\(s_2^a\)};
    \draw[->, dashed] (.55,1.01) to[bend left=25] (.88,.67);
  \end{tikzpicture}
  \caption{The world lines of isotropic observers in spacetime. By definition of
  isotropy, for any two vectors \(s_1^a,s_2^a\) at \(p\) which are orthogonal to
  \(u^a\), there exists an isometry of the spacetime which leaves \(p\) fixed and
  rotates \(s_1^a\) into \(s_2^a\).}
  \label{fig:5.2}
\end{figure}
\endgroup
